%% common/i18n/strings-fr.tex
%% Copyright 2025 Esther Poniatowski <esther.poniatowski@ens.psl.eu>
%
% This work may be distributed and/or modified under the
% conditions of the LaTeX Project Public License, either version 1.3c
% of this license or (at your option) any later version.
% The latest version of this license is in
%   https://www.latex-project.org/lppl.txt
% and version 1.3c or later is part of all distributions of LaTeX
% version 2008 or later.
%
% This work has the LPPL maintenance status `maintained'.
%
% The Current Maintainer of this work is Esther Poniatowski.
%
% This work consists of all files listed in MANIFEST.md.
%
%% Version: 2026/01/31 v0.1.0
% =============================================================================
% FRENCH LANGUAGE STRINGS (strings-fr.tex)
% =============================================================================
%
% PURPOSE:
%   Defines all user-facing text strings in French for the scholia package.
%   This enables internationalization (i18n) by abstracting hardcoded strings.
%
% USAGE:
%   Load this file in master.tex after packages but before other macros.
%   To create a new language, copy this file and translate all \def commands.
%
% NAMING CONVENTION:
%   \scholia@str@<system>@<identifier>
%   Example: \scholia@str@exo@label for "Exercice"
%
% =============================================================================

% -----------------------------------------------------------------------------
% EXERCISE SYSTEM STRINGS
% -----------------------------------------------------------------------------

% Main labels
\def\scholia@str@exo@label{Exercice}
\def\scholia@str@exo@guided{Guidé}
\def\scholia@str@exo@semiguided{Avec indications}
\def\scholia@str@exo@autonomous{Autonome}
\def\scholia@str@exo@modeling{Modélisation}
\def\scholia@str@exo@correction{Correction}

% Exercise sub-sections
\def\scholia@str@exo@hint{Indication.}
\def\scholia@str@exo@further{Pour aller plus loin.}

% -----------------------------------------------------------------------------
% GRADING SYSTEM STRINGS (BARÈME)
% -----------------------------------------------------------------------------

% Headers and labels
\def\scholia@str@bareme@label{Barème}
\def\scholia@str@bareme@criterion{Critère}
\def\scholia@str@bareme@description{Description}
\def\scholia@str@bareme@points{Points}
\def\scholia@str@bareme@total{Total}
\def\scholia@str@bareme@question{Question}

% Grading contexts
\def\scholia@str@bareme@exercise{Notation de l'exercice}
\def\scholia@str@bareme@totalgrade{Notation totale}
\def\scholia@str@bareme@competency{Compétence}
\def\scholia@str@bareme@criteria{Critères de réussite}
\def\scholia@str@bareme@resources{Ressources à consulter}

% Point suffix (e.g., "3 pts")
\def\scholia@str@bareme@pointsuffix{pts}
\def\scholia@str@bareme@pointsingular{pt}

% -----------------------------------------------------------------------------
% TYPED PAGE SYSTEM STRINGS
% -----------------------------------------------------------------------------

% Type display names
\def\scholia@str@typedpage@notion{Notions}
\def\scholia@str@typedpage@exercise{Exercices}
\def\scholia@str@typedpage@solution{Corrections}
\def\scholia@str@typedpage@autocheck{Auto-évaluation}
\def\scholia@str@typedpage@method{Méthode}
\def\scholia@str@typedpage@assessment{Évaluation}
\def\scholia@str@typedpage@summary{Sommaire}
\def\scholia@str@typedpage@keypoint{À retenir}
\def\scholia@str@typedpage@research{Recherche}
\def\scholia@str@typedpage@presentation{Présentation}
\def\scholia@str@typedpage@diagnostic{Fiche individuelle}

% Header elements
\def\scholia@str@typedpage@refprefix{Réf.:\;}
\def\scholia@str@typedpage@page{Page}
\def\scholia@str@typedpage@of{/}

% -----------------------------------------------------------------------------
% CALLOUT SYSTEM STRINGS
% -----------------------------------------------------------------------------

% Default titles (when no custom title provided)
\def\scholia@str@callout@attention{Attention}
\def\scholia@str@callout@definition{Définition}
\def\scholia@str@callout@hint{Astuce}
\def\scholia@str@callout@method{Méthode}
\def\scholia@str@callout@notations{Notations}
\def\scholia@str@callout@objective{Objectif}
\def\scholia@str@callout@proposition{Proposition}
\def\scholia@str@callout@proof{Démonstration}
\def\scholia@str@callout@remark{Remarque}
\def\scholia@str@callout@fail{Erreur fréquente}
\def\scholia@str@callout@instructions{Consignes}
\def\scholia@str@callout@summary{Bilan}
\def\scholia@str@callout@question{Question}
\def\scholia@str@callout@example{Exemple}
\def\scholia@str@callout@approach{Approche}

% -----------------------------------------------------------------------------
% SUMMARY TABLE STRINGS (SOMMAIRE)
% -----------------------------------------------------------------------------

\def\scholia@str@sommaire@competency{Compétence}
\def\scholia@str@sommaire@resources{Ressources à consulter}
\def\scholia@str@sommaire@reference{Référence}
\def\scholia@str@sommaire@title{Titre}

% -----------------------------------------------------------------------------
% PROGRESSION STRINGS (BEAMER)
% -----------------------------------------------------------------------------

\def\scholia@str@progression@session{Séance}

% -----------------------------------------------------------------------------
% NUMBERING STRINGS
% -----------------------------------------------------------------------------

\def\scholia@str@num@step{Étape}
\def\scholia@str@num@question{Question}

% -----------------------------------------------------------------------------
% METADATA STRINGS
% -----------------------------------------------------------------------------

\def\scholia@str@meta@chapter{Chapitre}
\def\scholia@str@meta@author{Auteur}
\def\scholia@str@meta@year{Année}
\def\scholia@str@meta@class{Classe}

% =============================================================================
% END OF FRENCH STRINGS
% =============================================================================
